\section*{Bloque A. Fundamentos teóricos y conceptuales}

\noindent\textit{Preguntas 1 a 30. Evalúan la demanda derivada de trabajo, el producto marginal, la oferta de trabajo, los efectos ingreso y sustitución, el monopsonio y el monopolio bilateral.}

\q{La demanda de trabajo es una demanda derivada porque:}
\begin{opts}
    \item depende únicamente de la oferta agregada de mano de obra.
    \item es fijada unilateralmente mediante negociación de sindicatos.
    \item depende de la demanda del bien final y de la productividad laboral.
    \item resulta completamente independiente de los salarios de mercado.
\end{opts}

\q{Una empresa competitiva maximizadora de beneficios contrata trabajo hasta que:}
\begin{opts}
    \item el ingreso marginal del producto iguala al salario.
    \item el salario real de mercado converge a cero soles.
    \item el producto marginal del trabajo alcanza su máximo.
    \item el costo marginal de producción cae hasta cero.
\end{opts}

\q{El ingreso marginal del producto del trabajo (MRP) es:}
\begin{opts}
    \item el salario nominal pagado por hora.
    \item el costo marginal de contratar trabajo.
    \item la productividad media física del factor.
    \item el producto marginal por el ingreso marginal.
\end{opts}

\q{La curva de demanda de trabajo de la empresa competitiva tiene pendiente negativa porque:}
\begin{opts}
    \item el producto marginal del trabajo cae al subir el empleo.
    \item el salario de mercado permanece estrictamente constante.
    \item la oferta agregada de trabajo es perfectamente vertical.
    \item el producto marginal laboral es constante a cualquier escala.
\end{opts}

\q{La demanda de trabajo es más elástica cuando:}
\begin{opts}
    \item el factor trabajo carece de sustitutos cercanos.
    \item el gasto laboral pesa poco en los costos totales.
    \item el bien final es elástico y hay sustitutos laborales.
    \item el salario nominal queda rígido por norma estatal.
\end{opts}

\q{Un aumento del salario genera:}
\begin{opts}
    \item únicamente un efecto escala de expansión.
    \item únicamente un efecto sustitución hacia capital.
    \item efectos escala y sustitución que reducen empleo.
    \item un desplazamiento expansivo de la demanda laboral.
\end{opts}

\q{La oferta de trabajo individual refleja:}
\begin{opts}
    \item la elección trabajo-ocio sujeta a tiempo e ingreso.
    \item únicamente las metas productivas de las empresas.
    \item la demanda agregada de bienes y servicios finales.
    \item el costo marginal de contratación en la industria.
\end{opts}

\q{Un aumento del salario produce un efecto sustitución que:}
\begin{opts}
    \item incentiva trabajar más y consumir menos ocio.
    \item incentiva trabajar menos y demandar más ocio.
    \item deja inalteradas las horas ofrecidas de trabajo.
    \item reduce necesariamente el ingreso laboral total.
\end{opts}

\q{La curva de oferta de trabajo puede tener pendiente negativa cuando:}
\begin{opts}
    \item el efecto sustitución supera al efecto ingreso.
    \item domina el efecto ingreso y se elige más ocio.
    \item el salario de mercado es exactamente igual a cero.
    \item la demanda de trabajo es infinitamente elástica.
\end{opts}

\q{El salario de reserva es:}
\begin{opts}
    \item el salario máximo que ofrece pagar una firma.
    \item el salario mínimo que induce a aceptar empleo.
    \item el salario competitivo de equilibrio de mercado.
    \item el costo marginal de contratación monopsonista.
\end{opts}

\q{La oferta de trabajo del mercado generalmente tiene pendiente:}
\begin{opts}
    \item vertical.
    \item negativa.
    \item positiva.
    \item horizontal.
\end{opts}

\q{La elasticidad de la oferta de trabajo mide:}
\begin{opts}
    \item la variación salarial provocada por cambios del empleo.
    \item la productividad física marginal del factor trabajo.
    \item la sensibilidad del trabajo ofrecido ante el salario.
    \item el costo marginal de expandir la planilla laboral.
\end{opts}

\q{Un monopsonio en el mercado de trabajo es:}
\begin{opts}
    \item muchos compradores de trabajo.
    \item un solo vendedor de trabajo.
    \item muchas empresas y un sindicato.
    \item un único comprador de trabajo.
\end{opts}

\q{Para el monopsonista, el costo marginal del factor trabajo es mayor que el salario porque:}
\begin{opts}
    \item el salario de mercado es perfectamente constante.
    \item contratar más exige subir el salario a toda la planta.
    \item la empresa monopsonista carece de costos marginales.
    \item la curva de oferta laboral es totalmente horizontal.
\end{opts}

\q{El monopsonista maximiza beneficios contratando hasta que:}
\begin{opts}
    \item el ingreso marginal del producto iguala al costo marginal factorial.
    \item el salario pagado iguala exactamente al costo marginal del factor.
    \item el producto físico marginal del trabajador se vuelve igual a cero.
    \item el salario nominal coincide exactamente con el salario mínimo legal.
\end{opts}

\q{En comparación con un mercado competitivo, el monopsonio:}
\begin{opts}
    \item paga un salario menor y contrata menos trabajo.
    \item paga un salario mayor y contrata más.
    \item paga el mismo salario y contrata lo mismo.
    \item no afecta el empleo.
\end{opts}

\q{La explotación monopsonística se refiere a:}
\begin{opts}
    \item un salario mayor que el MRP.
    \item un salario igual al MRP.
    \item un salario igual al costo marginal.
    \item un salario menor que el MRP.
\end{opts}

\q{Un salario mínimo puede aumentar el empleo en un monopsonio si:}
\begin{opts}
    \item se fija por debajo del salario de monopsonio.
    \item se fija entre el salario monopsonista y competitivo.
    \item se fija muy por encima del ingreso marginal producto.
    \item se prohíbe legalmente la negociación con sindicatos.
\end{opts}

\q{Un monopolio bilateral en el mercado de trabajo ocurre cuando:}
\begin{opts}
    \item existen numerosos compradores y vendedores atomizados.
    \item opera un único vendedor de trabajo frente a muchos compradores.
    \item un único comprador de trabajo negocia con un único sindicato.
    \item el Estado decreta directamente el salario de equilibrio.
\end{opts}

\q{En un monopolio bilateral, el salario y el empleo:}
\begin{opts}
    \item quedan fijados únicamente por oferta y demanda competitivas.
    \item igualan necesariamente el nivel de equilibrio competitivo.
    \item se reducen siempre a cero por conflicto distributivo.
    \item quedan indeterminados dentro de un rango negociable.
\end{opts}

\q{Un sindicato puede tener como objetivo:}
\begin{opts}
    \item elevar salarios, mejorar condiciones o priorizar empleo.
    \item minimizar el salario real pagado a los trabajadores.
    \item suprimir por completo el nivel de empleo contratado.
    \item fijar unilateralmente el precio final del producto.
\end{opts}

\q{La curva de contrato en la negociación salarial:}
\begin{opts}
    \item contiene pares eficientes de salario y empleo.
    \item coincide exactamente con la demanda de trabajo.
    \item representa la curva de oferta laboral individual.
    \item muestra el costo marginal salarial de la firma.
\end{opts}

\q{El poder de negociación de un sindicato depende de:}
\begin{opts}
    \item exclusivamente la elasticidad de demanda laboral.
    \item únicamente el número de empresas de la industria.
    \item sustituibilidad, normas, cohesión y costo de conflicto.
    \item solamente la tasa de interés real del mercado.
\end{opts}

\q{Una huelga es:}
\begin{opts}
    \item una forma de subsidio directo al empleo.
    \item un mecanismo de presión que impone costos.
    \item un impuesto específico sobre la planilla.
    \item una norma legal sobre el salario mínimo.
\end{opts}

\q{Según las reglas de Hicks-Marshall, la demanda de trabajo es más elástica cuando:}
\begin{opts}
    \item existen sustitutos, producto elástico y alto peso laboral.
    \item no existen sustitutos y el costo laboral pesa muy poco.
    \item la curva de oferta de trabajo es totalmente vertical.
    \item el salario de equilibrio queda fijado por decreto ley.
\end{opts}

\q{El cambio tecnológico que ahorra trabajo:}
\begin{opts}
    \item aumenta invariablemente la demanda de trabajo en todo sector.
    \item resulta completamente neutral sobre la contratación laboral.
    \item reduce empleo por sustitución, aunque puede subirlo por escala.
    \item suprime totalmente el mercado laboral en el largo plazo.
\end{opts}

\q{Si el trabajo y el capital son sustitutos, un aumento del salario:}
\begin{opts}
    \item aumenta directamente la demanda de trabajo.
    \item tiende a elevar la demanda de capital.
    \item reduce el precio de alquiler del capital.
    \item deja inalterada la demanda de capital.
\end{opts}

\q{Un diferencial salarial compensatorio es:}
\begin{opts}
    \item una discriminación de precios por grupo.
    \item un ajuste salarial por riesgo o condición.
    \item un subsidio estatal directo a la planilla.
    \item un salario mínimo uniforme intersectorial.
\end{opts}

\q{La inversión en capital humano:}
\begin{opts}
    \item reduce sistemáticamente la productividad laboral.
    \item carece de efecto alguno sobre los salarios reales.
    \item beneficia exclusivamente a los accionistas de firmas.
    \item eleva la productividad y los salarios esperados.
\end{opts}

\q{El equilibrio del mercado de trabajo competitivo se alcanza cuando:}
\begin{opts}
    \item la oferta total de trabajo cae hasta ser cero.
    \item la cantidad demandada iguala la cantidad ofrecida.
    \item el salario real de equilibrio converge a cero.
    \item el ingreso marginal del producto iguala a cero.
\end{opts}
